\section{Technical Overview}
\label{sec:technical-overview}

\textsc{Terrae} proves GBDT training by exposing the algebra already present in
histogram-based tree learning, rather than compiling the trainer into a large
generic circuit. 
In this section, we give a high-level overview of the main technical ingredients,
including how to batch claims across heterogeneous domains and how to prove histogram aggregation without a RAM argument.


\subsection{Batching Across Heterogeneous Domains}\label{subsec:overview-batching}

Direct algebraic certification creates many relations over many different
domains. Training uses tensors indexed by rounds, classes, samples, nodes,
features, and bins; lookup, range, prefix-sum, division, maximum, and
zero-check relations therefore have different lengths and shapes. Ordinary
same-domain batching would force extra commitments or leave many small claims
unbatched. \textsc{Terrae} introduces two batching tools for this setting.

\paragraph{Domain-lifting batching.}
A low-degree algebraic constraint over $\mathbb{H}_n$ is specified by a
constant-arity, constant-degree polynomial $\phi(x_1,\ldots,x_k)$. Vectors
$\mathbf{v}_1,\ldots,\mathbf{v}_k\in\mathbb{F}^n$ satisfy it when
\begin{align}
    \forall i \in [n].\
    \phi(\mathbf{v}_1[i],\ldots,\mathbf{v}_k[i])=0\\
    \iff \forall X\in \mathbb{H}_n.\ \phi(v_1(X),\ldots,v_k(X))=0.
\end{align}
If $v_1(X),\ldots,v_k(X)$ encode these vectors over $\mathbb{H}_n$ and are comitted, the polynomial
$f(X)=\phi(v_1(X),\ldots,v_k(X))$ is virtual: opening $f(t)$ reduces to
opening the $v_j(t)$ values and 
the verifier evaluating $\phi$ locally. For example,
$\phi(x,y)=x^2-y$ encodes the relation
$\mathbf{u}\odot\mathbf{u}=\mathbf{v}$.

Therefore every low-degree algebraic
constraint becomes a zero-check
for some possibly virtual $f(X)$:
\begin{align}
    \forall X\in \mathbb{H}_n.\ f(X)=0.
\end{align}
We prove this zero-check through the divisibility test
\begin{align}
    f(X)=q(X)\nu_n(X),
    \quad
    \nu_n(X)=X^n-1.
\end{align}
Here $\nu_n$ is the vanishing polynomial of $\mathbb{H}_n$.

Suppose two algebraic constraints live over nested power-of-two domains
$\mathbb{H}_{n_1}\subset \mathbb{H}_{n_2}$.
After converting them to zero-checks $f_1(X)$
over $\mathbb{H}_{n_1}$ and $f_2(X)$ over $\mathbb{H}_{n_2}$, instead of proving
them separately, \textsc{Terrae} lifts the smaller polynomial to the larger
domain:
\begin{align}
    \widehat{f}_1(X)
    :=
    f_1(X)\frac{\nu_{n_2}(X)}{\nu_{n_1}(X)}.
\end{align}
The verifier can then batch the lifted zero-check with the larger one as
\begin{align}
    \widehat{f}_1(X)+\eta f_2(X)=q(X)\nu_{n_2}(X)
\end{align}
for a random challenge $\eta$.

A bonus here is that the Aurora sum-check~\cite{aurora19}
also reduces to a divisibility test,
and hence we can batch sum-check claims
in a similar manner. See \Cref{thm:domain-lifting-batching} for details.

\paragraph{Interleaving batching.}
For a non-native predicate $\varphi$, a vector
$\mathbf{v}\in \mathbb{F}^n$ satisfies the predicate when
\begin{align}
    \forall i\in [n].\ \varphi(\mathbf{v}[i])=1.
\end{align}
Here $\varphi$ is not a low-degree field polynomial, so the verifier cannot
derive a virtual commitment as above. Examples include table
membership, range, and fixed-point division.

We batch such constraints by
interleaving same-predicate instances. For two vectors
$\mathbf{f},\mathbf{g}\in\mathbb{F}^n$ satisfying the same predicate $\phi$,
\textsc{Terrae} forms the virtual interleaving
\begin{align}
    (\mathbf{f}[0],\mathbf{g}[0],\mathbf{f}[1],\mathbf{g}[1],\ldots).
\end{align}
In polynomial form, if $f$ and $g$ are encoded over $\mathbb{H}_n$, the merged
object over $\mathbb{H}_{2n}$ is
\begin{align}
    \operatorname{Int}_n(f,g)(X)
    :=
    \operatorname{even}_n(X)f(X)
    +\operatorname{odd}_n(X)g(\omega_{2n}^{-1}X),
\end{align}
where
\begin{align}
    \operatorname{even}_n(X)=\frac{X^n+1}{2},
    \qquad
    \operatorname{odd}_n(X)=\frac{1-X^n}{2}.
\end{align}
Note that both $\operatorname{even}_n(X)$ and $\operatorname{odd}_n(X)$ are sparse, and the verifier
can evaluate them at any point with $O(\log n)$ field operations.

For vector $\mathbf{f}\in \mathbb{F}^n$ and $\mathbf{g}\in \mathbb{F}^{nr}$ of different lengths, we first lift the shorter vector to the larger domain by a
substitution $X\mapsto X^r$, i.e. defining
\begin{align}
    \widehat{g}(X):=g(X^r).
\end{align}
As stated in \Cref{lem:row-replication}, $\widehat{\mathbf{g}}$ replicates $\mathbf{g}$ over domain $\mathbb{H}_{nr}$ for $r$ times. 
We then interleave $\mathbf{f}$ and $\widehat{\mathbf{g}}$ 
and derive the final interleaved polynomial,
which is still virtual: 
its openings are derived from
evaluations of existing commitments and public sparse selectors, so batching
does not require committing to new auxiliary witness polynomials.

We present the formal interleaving batching method
in \Cref{lem:interleaving}, and use it to batch all non-native predicates in \textsc{Terrae}.

\subsection{Histogram Aggregation Without RAM}\label{subsec:overview-histogram}

GBDT training repeatedly converts sample-wise information into per-node,
per-feature, and per-bin statistics. Prior certification-style systems~\cite{sparrow24,zkboost26}
naturally model this as a memory update process: for each routed sample, update
the corresponding histogram bucket and prove the resulting RAM trace.
\textsc{Terrae} instead uses a histogram-specific argument.

For simplicity, consider a simplified version of histogram aggregation.
Suppose we have vectors $\mathbf{ind}\in [m]^n$,
$\mathbf{u}\in \mathbb{F}^n$, and $\mathbf{v}\in \mathbb{F}^m$.
Our goal is to prove that $\mathbf{v}$ is the histogram
of $\mathbf{u}$ according to the indices in $\mathbf{ind}$:
\begin{align}
    \mathbf{v}[j]=\sum_{i:\mathbf{ind}[i]=j}\mathbf{u}[i].
\end{align}

To prove this, the verifier generates a random challenge
$\alpha\gets \mathbb{F}$, which binds each $i$-th slot to $\alpha^i$. The parties next move to prove
\begin{align}
    \sum_{i=0}^{n-1}\mathbf{u}[i]\cdot \alpha^{\mathbf{ind}[i]}
    =
    \sum_{j=0}^{m-1}\mathbf{v}[j]\cdot \alpha^j.\label{eq:histogram-aggregation-power-sum}
\end{align}
The remaining question is how to prove the exponentiated terms
$\alpha^{\mathbf{ind}[i]}$.
The prover commits a power table
$\mathbf{p}\in\mathbb{F}^{m}$ with intended entries
$\mathbf{p}[j]=\alpha^j$. It proves that $\mathbf{p}$ is well formed by
checking
\begin{align}
    \mathbf{p}[0]&=1,\\
    \mathbf{p}[j]&=\alpha\cdot \mathbf{p}[j-1],
    \qquad j\in\{1,\ldots,m-1\}.
\end{align}
Or equivalently in polynomial form,
\begin{align}
    &p(\omega_m^0) = 1,\\
    &(p(t) - \alpha\cdot p(\omega_m^{-1} t))\cdot (1 - \lambda_m^{(0)}(t)) = q(X)\nu_{m}(X),
\end{align}
where recall $\lambda_m^{(0)}(X)$ is a Lagrange
basis polynomial, and the second relation 
is a zero-check.

The prover then proves by lookup that each row
$(\mathbf{ind}[i],\mathbf{pow}[i])$ appears in the public-indexed table
$(j,\mathbf{pow}[j])_{j\in[m]}$, i.e.
\begin{align}
    \mathbf{ind}\| \mathbf{pow}\subseteq (0, 1, \dots, m-1)^\top\| \mathbf{p}.
\end{align}
This binds
$\mathbf{pow}[i]=\alpha^{\mathbf{ind}[i]}$, after which
\Cref{eq:histogram-aggregation-power-sum}
becomes
\begin{align}
    \sum_{i=0}^{n-1}\mathbf{u}[i]\cdot \mathbf{pow}[i]
    =
    \sum_{j=0}^{m-1}\mathbf{v}[j]\cdot \mathbf{p}[j],
\end{align}
which can be checked by standard sum-check.

Soundness follows from the randomness of $\alpha$. 
To analyze the soundness error introduced by
\Cref{eq:histogram-aggregation-power-sum},
define
\[
    \Delta_j
    =
    \mathbf{v}[j]
    -
    \sum_{i:\mathbf{ind}[i]=j}\mathbf{u}[i].
\]
If the claim that $\mathbf{v}$ is the histogram of $\mathbf{u}$ is false, 
then some $\Delta_j\neq 0$, and the accepted aggregate identity implies
\[
    \sum_{j=0}^{m-1}\Delta_j\alpha^j=0.
\]
This is a nonzero polynomial in $\alpha$ of degree at most $m-1$, so a cheating
prover can satisfy it with probability at most $(m-1)/|\mathbb{F}|$, assuming
$\alpha$ is sampled after the relevant commitments are fixed. Thus the
aggregate identity is a compact randomized check of the full histogram relation.
